%% Created by Maple 2022.1, Windows 10
%% Source Worksheet: lab_sols01
%% Generated: Wed May 08 10:37:48 BST 2024
\documentclass{article}
\usepackage{amssymb}
\usepackage{graphicx}
\usepackage{hyperref}
\usepackage{listings}
\usepackage{mathtools}
\usepackage{maple}
\usepackage[utf8]{inputenc}
\usepackage[svgnames]{xcolor}
\usepackage{amsmath}
\usepackage{breqn}
\usepackage{textcomp}
\begin{document}
\lstset{basicstyle=\ttfamily,breaklines=true,columns=flexible}
\pagestyle{empty}
\DefineParaStyle{Maple Bullet Item}
\DefineParaStyle{Maple Heading 1}
\DefineParaStyle{Maple Warning}
\DefineParaStyle{Maple Heading 4}
\DefineParaStyle{Maple Heading 2}
\DefineParaStyle{Maple Heading 3}
\DefineParaStyle{Maple Dash Item}
\DefineParaStyle{Maple Error}
\DefineParaStyle{Maple Title}
\DefineParaStyle{Maple Ordered List 1}
\DefineParaStyle{Maple Text Output}
\DefineParaStyle{Maple Ordered List 2}
\DefineParaStyle{Maple Ordered List 3}
\DefineParaStyle{Maple Normal}
\DefineParaStyle{Maple Ordered List 4}
\DefineParaStyle{Maple Ordered List 5}
\DefineCharStyle{Maple 2D Output}
\DefineCharStyle{Maple 2D Input}
\DefineCharStyle{Maple Maple Input}
\DefineCharStyle{Maple 2D Math}
\DefineCharStyle{Maple Hyperlink}
\begin{center}
\begin{Maple Title}
\underline{\textbf{Introduction to Maple}}
\end{Maple Title}
\end{center}
\begin{Maple Normal}
\_\_\_\_\_\_\_\_\_\_\_\_\_\_\_\_\_\_\_\_\_\_\_\_\_\_\_\_\_\_\_\_\_\_\_\_\_\_\_\_\_\_\_\_\_\_\_\_\_\_\_\_\_\_\_\_\_\_\_\_\_\_\_\_\_\_\_\_\_\_\_\_\_\_\_\_\_\_\_\_\_
\end{Maple Normal}
\begin{Maple Heading 2}
\textbf{Exercise 2.1}
\end{Maple Heading 2}
\begin{lstlisting}
> 2+2;
\end{lstlisting}
\begin{lstlisting}
> (3+7+10)*(1000-8)/(900+90+2) - 17;
\end{lstlisting}
\begin{Maple Normal}
\_\_\_\_\_\_\_\_\_\_\_\_\_\_\_\_\_\_\_\_\_\_\_\_\_\_\_\_\_\_\_\_\_\_\_\_\_\_\_\_\_\_\_\_\_\_\_\_\_\_\_\_\_\_\_\_\_\_\_\_\_\_\_\_\_\_\_\_\_\_\_\_\_\_\_\_\_\_\_\_\_
\end{Maple Normal}
\begin{Maple Heading 2}
\textbf{Exercise 2.2}
\end{Maple Heading 2}
\begin{Maple Normal}
\textbf{(a)}
\end{Maple Normal}
\begin{lstlisting}
> 6^20 * 15^20 / 9^20;
\end{lstlisting}
\begin{Maple Normal}
The answer is one followed by 20 zeros, or in other words 
$ 10^{20} $.  This is easy to see by hand, because 
$ \frac{6^{20} 15^{20}}{9^{20}}=(\frac{\mathit{6\,x\,15\,}}{9})^{20} $ and (6 x 15)/9 is just 10.
\end{Maple Normal}
\begin{Maple Normal}
\textbf{(b)}
\end{Maple Normal}
\begin{lstlisting}
> (10^10 - 1)/99;
\end{lstlisting}
\begin{Maple Normal}
The answer is 
$ 1+100+10000+1000000+100000000 $, or in other words 
$ 1+100+100^{2}+100^{3}+100^{4} $.  The standard geometric progression formula says that this is the same as 
$ \frac{100^{5}-1}{100-1} $, or in other words 
$ \frac{10^{10}-1}{99} $.
\end{Maple Normal}
\begin{Maple Normal}
\textbf{(c)}
\end{Maple Normal}
\begin{lstlisting}
> (10^10 - 10 - 9^2)/9^2;
\end{lstlisting}
\begin{Maple Normal}
To check this by hand, let 
$ x  $ be the number 
$ 123456789 $.  Then 
$ 10 x =1234567890 $ and so 
$ 10 x +9=1234567899 $.  If we subtract 
$ x  $ from this the digits mostly cancel and we get 
$ 9 x +9=1111111110 $.  Multiply by 9 again to get 
$ 9^{2} (x +1)=9999999990 $, which is 
$ 10^{10}-10 $.  Rearrange this to get 
$ x +1=\frac{10^{10}-10}{9^{2}} $ and so 
$ x =\frac{10^{10}-10-9^{2}}{9^{2}} $.
\end{Maple Normal}
\begin{Maple Normal}
\textbf{(d)}
\end{Maple Normal}
\begin{lstlisting}
> (10^9+1)*(10^10 - 10 - 9^2)/9^2;
\end{lstlisting}
\begin{Maple Normal}
Assuming part (c), we have 
$ \frac{10^{10}-10-9^{2}}{9^{2}}=123456789 $ and so 
$ \frac{10^{9} (10^{10}-10-9^{2})}{9^{2}}=123456789000000000 $.  These two numbers can be added without any carrying, to give
$ \frac{(10^{9}+1) (10^{10}-10-9^{2})}{9^{2}}=123456789000000000+123456789 $ = 123456789123456789.
\end{Maple Normal}
\begin{Maple Normal}
\_\_\_\_\_\_\_\_\_\_\_\_\_\_\_\_\_\_\_\_\_\_\_\_\_\_\_\_\_\_\_\_\_\_\_\_\_\_\_\_\_\_\_\_\_\_\_\_\_\_\_\_\_\_\_\_\_\_\_\_\_\_\_\_\_\_\_\_\_\_\_\_\_\_\_\_\_\_\_\_\_
\end{Maple Normal}
\begin{Maple Heading 2}
\textbf{Exercise 2.3}
\end{Maple Heading 2}
\begin{Maple Normal}
The numbers in this exercise are approximations to 
$ \mathrm{Pi} $; in a certain sense, they are actually the best possible approximations.
\end{Maple Normal}
\begin{lstlisting}
> 3 + 1/7;
\end{lstlisting}
\begin{lstlisting}
> evalf(%);
\end{lstlisting}
\begin{lstlisting}
> 3 + 1/(7 + 1/15);
\end{lstlisting}
\begin{lstlisting}
> evalf(%);
\end{lstlisting}
\begin{lstlisting}
> 3+1/(7+1/(15 + 1/(1 + 1/293)));
\end{lstlisting}
\begin{lstlisting}
> evalf(%);
\end{lstlisting}
\begin{lstlisting}
> evalf(Pi);
\end{lstlisting}
\begin{lstlisting}
> 
\end{lstlisting}
\begin{Maple Normal}
\_\_\_\_\_\_\_\_\_\_\_\_\_\_\_\_\_\_\_\_\_\_\_\_\_\_\_\_\_\_\_\_\_\_\_\_\_\_\_\_\_\_\_\_\_\_\_\_\_\_\_\_\_\_\_\_\_\_\_\_\_\_\_\_\_\_\_\_\_\_\_\_\_\_\_\_\_\_\_\_\_
\end{Maple Normal}
\begin{Maple Heading 2}
\textbf{Exercise 2.4}
\end{Maple Heading 2}
\begin{lstlisting}
> evalf[4997](exp(1));
\end{lstlisting}
\begin{Maple Normal}

\end{Maple Normal}
\begin{Maple Normal}
\_\_\_\_\_\_\_\_\_\_\_\_\_\_\_\_\_\_\_\_\_\_\_\_\_\_\_\_\_\_\_\_\_\_\_\_\_\_\_\_\_\_\_\_\_\_\_\_\_\_\_\_\_\_\_\_\_\_\_\_\_\_\_\_\_\_\_\_\_\_\_\_\_\_\_\_\_\_\_\_\_
\end{Maple Normal}
\begin{Maple Heading 2}
\textbf{Exercise 2.5}
\end{Maple Heading 2}
\begin{lstlisting}
> 
\end{lstlisting}
\begin{lstlisting}
> Digits := 40;
\end{lstlisting}
\begin{lstlisting}
> x := exp(Pi*sqrt(163));
\end{lstlisting}
\begin{lstlisting}
> y := evalf(x);
\end{lstlisting}
\begin{lstlisting}
> printf("%20.20f\n\n",y);
\end{lstlisting}
\begin{lstlisting}
> z := round(y);
\end{lstlisting}
\begin{lstlisting}
> y-z;
\end{lstlisting}
\begin{lstlisting}
> restart;
\end{lstlisting}
\begin{lstlisting}
> 
\end{lstlisting}
\begin{Maple Normal}
\_\_\_\_\_\_\_\_\_\_\_\_\_\_\_\_\_\_\_\_\_\_\_\_\_\_\_\_\_\_\_\_\_\_\_\_\_\_\_\_\_\_\_\_\_\_\_\_\_\_\_\_\_\_\_\_\_\_\_\_\_\_\_\_\_\_\_\_\_\_\_\_\_\_\_\_\_\_\_\_\_
\end{Maple Normal}
\begin{Maple Heading 2}
\textbf{Exercise 3.1}
\end{Maple Heading 2}
\begin{lstlisting}
> A := (x^2-4*y^2)*(x^3-x*y^2);
\end{lstlisting}
\begin{lstlisting}
> simplify(A);
\end{lstlisting}
\begin{lstlisting}
> expand(A);
\end{lstlisting}
\begin{lstlisting}
> factor(A);
\end{lstlisting}
\begin{lstlisting}
> convert(expand(A),horner,x);
\end{lstlisting}
\begin{Maple Normal}
I would vote for 
$ (x -2 y) (x +2 y) x (x -y) (x +y) $ as the most useful version, but of course it depends what you want to use it for.
\end{Maple Normal}
\begin{Maple Normal}
\_\_\_\_\_\_\_\_\_\_\_\_\_\_\_\_\_\_\_\_\_\_\_\_\_\_\_\_\_\_\_\_\_\_\_\_\_\_\_\_\_\_\_\_\_\_\_\_\_\_\_\_\_\_\_\_\_\_\_\_\_\_\_\_\_\_\_\_\_\_\_\_\_\_\_\_\_\_\_\_\_
\end{Maple Normal}
\begin{Maple Heading 2}
\textbf{Exercise 3.2}
\end{Maple Heading 2}
\begin{lstlisting}
> simplify(2*x/(x^2-1)+1/(x+x^2)+1/(x-x^2));
\end{lstlisting}
\begin{Maple Normal}
To do this by hand, note that 
$ x^{2}-1=(x +1)(x -1) $ and 
$ x +x^{2}=x (x +1) $ and 
$ x -x^{2}=-x (x -1) $.  Thus, our expression is 
$ \frac{2 x^{2}}{x (x +1) (x -1)}+\frac{x -1}{x (x +1) (x -1)}-\frac{x +1}{x (x +1) (x -1)} $, which simplifies to 
$ \frac{2 x^{2}-2}{x (x +1) (x -1)} $.  The numerator here factors as 
$ 2 (x +1) (x -1) $ and so everything cancels to leave 
$ \frac{2}{x} $.
\end{Maple Normal}
\begin{lstlisting}
> 
\end{lstlisting}
\begin{Maple Normal}
\_\_\_\_\_\_\_\_\_\_\_\_\_\_\_\_\_\_\_\_\_\_\_\_\_\_\_\_\_\_\_\_\_\_\_\_\_\_\_\_\_\_\_\_\_\_\_\_\_\_\_\_\_\_\_\_\_\_\_\_\_\_\_\_\_\_\_\_\_\_\_\_\_\_\_\_\_\_\_\_\_
\end{Maple Normal}
\begin{Maple Heading 2}
\textbf{Exercise 4.1}
\end{Maple Heading 2}
\begin{lstlisting}
> plot(sin(x),x=-5..5);
\end{lstlisting}
\begin{lstlisting}
> plot(x-x^3/6+x^5/120-x^7/5040,x=-5..5);
\end{lstlisting}
\begin{lstlisting}
> plot([sin(x),x-x^3/6+x^5/120-x^7/5040],x=-5..5);
\end{lstlisting}
\begin{Maple Normal}
The two graphs are very close together for 
$ x  $ between about 
$ -3 $ and 
$ 3 $, but outside that range they move apart very rapidly.
\end{Maple Normal}
\begin{Maple Normal}
\_\_\_\_\_\_\_\_\_\_\_\_\_\_\_\_\_\_\_\_\_\_\_\_\_\_\_\_\_\_\_\_\_\_\_\_\_\_\_\_\_\_\_\_\_\_\_\_\_\_\_\_\_\_\_\_\_\_\_\_\_\_\_\_\_\_\_\_\_\_\_\_\_\_\_\_\_\_\_\_\_
\end{Maple Normal}
\begin{Maple Heading 2}
\textbf{Exercise 5.1}
\end{Maple Heading 2}
\begin{lstlisting}
>   solve(\{x^2+y^2=1,(x-1)^2+(y-1)^2=1\},\{x,y\});
\end{lstlisting}
\begin{Maple Normal}
To obtain this by hand, expand out the second equation to get 
$ x^{2}-2 x +1+y^{2}-2 y +1=1 $, or equivalently 
$ x^{2}+y^{2}=2 x +2 y -1 $.  Subtracting the first equation gives 
$ 0=2 x +2 y -2 $, so 
$ x +y =1 $.  Squaring this gives 
$ x^{2}+y^{2}+2 x y =1 $, and subtracting the first equation gives 
$ 2 x y =0 $.  This means that either 
$ x  $ or 
$ y  $ must be zero, and the equation 
$ x +y =1 $ means that the other one must be 
$ 1 $.
\end{Maple Normal}
\begin{Maple Normal}

\end{Maple Normal}
\begin{Maple Normal}
Geometrically, the equation 
$ x^{2}+y^{2}=1 $ describes the circle of radius one centred at the origin, and the equation 
$ (x -1)^{2}+(y -1)^{2}=1 $ describes the circle of radius one centred at (1,1).  The two circles intersect at the points (1,0) and (0,1), which give the two solutions to our equations.
\end{Maple Normal}
\begin{Maple Normal}

\end{Maple Normal}
\begin{lstlisting}
> plots[display](
  plot([cos(t),sin(t),t=0..2*Pi],color=blue),
  plot([1+cos(t),1+sin(t),t=0..2*Pi],color=red)
);
\end{lstlisting}
\begin{lstlisting}
> 
\end{lstlisting}
\begin{Maple Normal}
\_\_\_\_\_\_\_\_\_\_\_\_\_\_\_\_\_\_\_\_\_\_\_\_\_\_\_\_\_\_\_\_\_\_\_\_\_\_\_\_\_\_\_\_\_\_\_\_\_\_\_\_\_\_\_\_\_\_\_\_\_\_\_\_\_\_\_\_\_\_\_\_\_\_\_\_\_\_\_\_\_
\end{Maple Normal}
\begin{Maple Heading 2}
\textbf{Exercise 6.1}
\end{Maple Heading 2}
\begin{lstlisting}
> diff(ln(ln(ln(x))),x);
\end{lstlisting}
\begin{lstlisting}
> (3*x+4)/(2*x+3);
\end{lstlisting}
\begin{lstlisting}
> diff(%,x);
\end{lstlisting}
\begin{lstlisting}
> simplify(%);
\end{lstlisting}
\begin{lstlisting}
> (1+x^2+x^4/2)*exp(-x^2);
\end{lstlisting}
\begin{lstlisting}
> diff(%,x);
\end{lstlisting}
\begin{lstlisting}
> simplify(%);
\end{lstlisting}
\begin{lstlisting}
> 
\end{lstlisting}
\begin{Maple Normal}
\_\_\_\_\_\_\_\_\_\_\_\_\_\_\_\_\_\_\_\_\_\_\_\_\_\_\_\_\_\_\_\_\_\_\_\_\_\_\_\_\_\_\_\_\_\_\_\_\_\_\_\_\_\_\_\_\_\_\_\_\_\_\_\_\_\_\_\_\_\_\_\_\_\_\_\_\_\_\_\_\_
\end{Maple Normal}
\begin{Maple Heading 2}
\textbf{Exercise 6.2}
\end{Maple Heading 2}
\begin{lstlisting}
> int(sqrt(1-x^2),x=0..1);
\end{lstlisting}
\begin{lstlisting}
> 
\end{lstlisting}
\end{document}
